\begin{lemma}
\label{lemma-ext-1-annihilated}
Let $R$ be ring and let $I \subset R$ be an ideal.
Let $K \in D(R)$ with $H^i(K) = 0$ for $i \not \in \{-1, 0\}$.
The following are equivalent
\begin{enumerate}
\item there exists a $c \geq 0$ such that the equivalent
conditions (1), (2), (3) of Lemma \ref{lemma-ext-1-annihilated-definite}
hold for $K$ and the ideal $I^c$,
\item there exists a $c \geq 0$ such that (a) $I^c$ annihilates
$H^{-1}(K)$ and (b) $H^0(K)$ is an $I^c$-projective module (see
Section \ref{section-near-projective}).
\end{enumerate}
If $R$ is Noetherian and $H^i(K)$ is a finite $R$-module
for $i = -1, 0$, then these are also equivalent to
\begin{enumerate}
\item[(3)] there exists a $c \geq 0$ such that the equivalent
conditions (4), (5) of Lemma \ref{lemma-ext-1-annihilated-definite}
hold for $K$ and the ideal $I^c$,
\item[(4)] $H^{-1}(K)$ is $I$-power torsion and there exist
$f_1, \ldots, f_s \in R$ with $V(f_1, \ldots, f_s) \subset V(I)$
such that the localizations $H^0(K)_{f_i}$ are projective
$R_{f_i}$-modules,
\item[(5)] $H^{-1}(K)$ is $I$-power torsion and there exist
$f_1, \ldots, f_s \in I$ with $V(f_1, \ldots, f_s) = V(I)$
such that the localizations $H^0(K)_{f_i}$ are projective
$R_{f_i}$-modules, and
\item[(6)] $H^{-1}(K)$ is $I$-power torsion and for any
$f_1, \ldots, f_s \in I$ with $V(f_1, \ldots, f_s) = V(I)$
the localizations $H^0(K)_{f_i}$ are projective $R_{f_i}$-modules.
\end{enumerate}
\end{lemma}

\begin{proof}
The distinguished triangle $H^{-1}(K)[1] \to K \to H^0(K)[0] \to H^{-1}(K)[2]$
determines an exact sequence
$$
0 \to \Ext^1_R(H^0(K), N) \to \Ext^1_R(K, N) \to \Hom_R(H^{-1}(K), N) \to
\Ext^2_R(H^0(K), N)
$$
Thus (2) implies that $I^{2c}$ annihilates $\Ext^1_R(K, N)$ for every
$R$-module $N$. Assuming (1) we immediately see that $H^0(K)$ is
$I^c$-projective. On the other hand, we may choose an injective map
$H^{-1}(K) \to N$ for some injective $R$-module $N$. Then this map
is the image of an element of $\Ext^1_R(K, N)$ by the vanishing
of the $\Ext^2$ in the sequence and we conclude $H^{-1}(K)$
is annihilated by $I^c$.

\medskip\noindent
Assume $R$ is Noetherian and $H^i(K)$ is a finite $R$-module
for $i = -1, 0$. By Lemma \ref{lemma-ext-1-annihilated-definite}
we see that (3) is equivalent to (1) and (2). Also, if (3)
holds then for $f \in I$ the multiplication by $f$ on
$H^0(K)$ factors through a projective module, which implies
that $H^0(K)_f$ is a summand of a projective $R_f$-module and hence
itself a projective $R_f$-module. Thus the equivalent conditions
(1), (2), and (3) imply (6).
Of course (6) implies (5) and (5) trivially implies (4).

\medskip\noindent
Assume (4). Since $H^{-1}(K)$ is a finite $R$-module and $I$-power torsion
we see that $I^{c_1}$ annihilates $H^{-1}(K)$ for some $c_1 \geq 0$.
Choose a short exact sequence
$$
0 \to M \to R^{\oplus r} \to H^0(K) \to 0
$$
which determines an element $\xi \in \Ext^1_R(H^0(K), M)$.
For any $f \in I$ we have $\Ext^1_R(H^0(K), M)_f = \Ext^1_{R_f}(H^0(K)_f, M_f)$
by Lemma \ref{lemma-pseudo-coherence-and-base-change-ext}.
Hence if $H^0(K)_f$ is projective, then a power of $f$ annihilates $\xi$.
We conclude that $\xi$ is annihilated by $(f_1, \ldots, f_s)^{c_2}$
for some $c_2 \geq 0$. Since $V(f_1, \ldots, f_s) \subset V(I)$ we have
$\sqrt{I} \subset (f_1, \ldots, f_s)$
(Algebra, Lemma \ref{algebra-lemma-Zariski-topology}).
Since $R$ is Noetherian we find $I^{c_3} \subset (f_1, \ldots, f_s)$
for some $c_3 \geq 0$ (Algebra, Lemma \ref{algebra-lemma-Noetherian-power}).
Hence $I^{c2c3}$ annihilates $\xi$.
This in turn says that $H^0(K)$ is $I^{c_2c_3}$-projective (as multiplication
by $a \in I$ which annihilate $\xi$ factor through $R^{\oplus r}$).
Hence taking $c = \max(c_1, c_2c_3)$ we see that (2) holds.
\end{proof}